\documentclass[12pt,a4paper]{article}
\usepackage[margin=25mm, headheight=15pt, headsep=10pt]{geometry}
\usepackage{fontspec}
\usepackage[table]{xcolor} % Note: table option is required for rowcolors
\usepackage{titlesec}
\usepackage{tocloft}
\usepackage{fancyhdr}
\usepackage{booktabs}
\usepackage{tcolorbox}
\tcbuselibrary{skins, breakable}
\usepackage[colorlinks=true, linkcolor=emerald, urlcolor=emerald, bookmarks=true]{hyperref}

\definecolor{emerald}{RGB}{0,102,51}
\definecolor{darkred}{RGB}{139,0,0}
\definecolor{lightgray}{RGB}{245,245,245}

\usepackage[nil]{babel}
\babelprovide[import, main]{bengali}
\babelprovide[import]{arabic}
\babelfont{rm}[Renderer=HarfBuzz,Scale=1.0]{Noto Serif Bengali}
\babelfont{sf}[Renderer=HarfBuzz]{Noto Sans Bengali}
\babelfont{tt}[Renderer=HarfBuzz]{Noto Sans Bengali}
\babelfont[arabic]{rm}[Renderer=HarfBuzz,Scale=1.1]{Noto Naskh Arabic}

% Section styling
\titleformat{\section}{\Large\bfseries\color{emerald}}{\thesection.}{0.5em}{}
\titleformat{\subsection}{\large\bfseries\color{emerald!85!black}}{\thesubsection.}{0.5em}{}
\titleformat{\subsubsection}{\normalsize\bfseries\color{emerald!70!black}}{\thesubsubsection.}{0.5em}{}
\titlespacing*{\section}{0pt}{18pt}{8pt}

% Fancy Header/Footer setup
\pagestyle{fancy}
\fancyhf{} % clear all header and footer fields
\fancyhead[L]{\color{emerald}\small\textbf{\nouppercase{\leftmark}}}
\fancyfoot[L]{\scriptsize\color{black!50}{\lat{AI}}-সংকলিত, আলেম-যাচাইকৃত নয়}
\fancyfoot[C]{\color{emerald!70!black}\thepage}
\renewcommand{\headrulewidth}{0.4pt}
\renewcommand{\headrule}{\hbox to\headwidth{\color{emerald}\leaders\hrule height \headrulewidth\hfill}}
\renewcommand{\footrulewidth}{0pt}

% Custom boxes
% 1. Ayat/Hadith Box
\newtcolorbox{ayatbox}[1][]{
    enhanced, breakable, 
    colback=emerald!5, 
    colframe=emerald, 
    boxrule=0.5pt, 
    arc=3pt, 
    left=10pt, right=10pt, top=10pt, bottom=10pt,
    #1
}

% 2. Quote/Translation Box
\newtcolorbox{quotebox}[1][]{
    enhanced, breakable,
    frame hidden,
    colback=lightgray,
    borderline west={3pt}{0pt}{emerald},
    left=15pt, right=10pt, top=10pt, bottom=10pt,
    sharp corners,
    #1
}

% 3. AI-generated notice box
\newtcolorbox{warnbox}[1][]{
    enhanced, breakable,
    colback=red!4,
    colframe=darkred,
    boxrule=0.8pt,
    arc=3pt,
    left=10pt, right=10pt, top=10pt, bottom=10pt,
    #1
}

% Commands
\newcommand{\arab}[1]{\foreignlanguage{arabic}{#1}}
\newcommand{\kib}[1]{\textcolor{emerald}{\textbf{#1}}}

% Latin (English) fallback: Noto Serif Bengali has NO Latin glyphs
\newfontfamily\latfont[Renderer=HarfBuzz]{Noto Serif}
\newcommand{\lat}[1]{{\latfont #1}}

\setlength{\parskip}{5pt}
\setlength{\parindent}{0pt}

% Fix for missing bullet in Noto Serif Bengali
\renewcommand{\labelitemi}{$\bullet$}



\begin{document}
\begin{center}{\Large\bfseries\color{emerald} 05-উমার-নবীজির-ঘরে-পরকালচিন্তা}\end{center}
\vspace{1em}
\section*{ঘটনা ৫: উমার (রাঃ) — নবীজির ঘরে পরকালের ভাবনা}

\begin{warnbox}
 \textbf{গুরুত্বপূর্ণ নোটিশ (\lat{Important} \lat{Notice}):} এই কন্টেন্টটি \lat{AI} (কৃত্রিম বুদ্ধিমত্তা) দিয়ে প্রস্তুত ও সংকলিত। \lat{AI} কঠোর সূত্র-মান নীতি মেনে শুধু নির্ভরযোগ্য উৎস থেকে তথ্য সংগ্রহ করেছে — কেবল সহীহ/হাসান হাদিস ও সহীহ সনদের আসার রাখা হয়েছে, দুর্বল সনদ বাদ দেওয়া হয়েছে। তবে এটি কোনো আলেম বা মুহাদ্দিস কর্তৃক যাচাইকৃত নয়।
\end{warnbox}


\medskip\hrule\medskip

\subsection*{১. সূত্র কার্ড}

\begin{table}[h]\centering\small
\begin{tabular}{ll}
বিষয় & বিবরণ \\
\hline
\textbf{বর্ণনাকারী} & উমার ইবনুল খাত্তাব (রাঃ) নিজেই; আবদুল্লাহ ইবনে আব্বাস (রাঃ) কর্তৃক বর্ণিত \\
\textbf{গ্রন্থ} & সুনানে ইবনে মাজাহ, হাদিস ৪১৫৩; সহীহ বুখারি, হাদিস ৪৯১৮ ও ২৪৬৩ (দীর্ঘ বর্ণনা) \\
\textbf{অধ্যায়} & কিতাবুজ-যুহদ, বাব: আলু মুহাম্মাদের বিছানা \\
\textbf{সনদের মান} & \textbf{সহীহ} (আলি জাই ও বুখারির সূত্রে) \\
\textbf{মূল বিষয়} & দুনিয়ার বিলাস বনাম আখেরাত — নবীজির ত্যাগে উমারের কান্না \\
\hline
\end{tabular}
\end{table}

\medskip\hrule\medskip

\subsection*{২. পটভূমি (কে, কখন, কোথায়, কেন)}

এটি সেই সময়কার ঘটনা, যখন নবীজি (সাল্লাল্লাহু আলাইহি ওয়া সাল্লাম) তাঁর স্ত্রীদের থেকে আলাদা হয়ে একটি উঁচু কামরায় অবস্থান করছিলেন — স্ত্রীদের সাথে নবীজির একটি \lat{disagreement} (মতপার্থক্য)-এর কারণে। উমার (রাঃ) \lat{news} (খবর) পেয়ে নবীজির কাছে এলেন।

সে সময় \textbf{কিসরা (পারস্যের সম্রাট) ও কায়সার (রোমের সম্রাট)} — বিশ্বের দুই পরাশক্তির শাসক — বিলাসী প্রাসাদে স্বর্ণ-রৌপ্য, ফল-ফুলের বাগানে \lat{luxury} (বিলাস) নিয়ে বসবাস করছিলেন। আর আল্লাহর নবী শুয়ে ছিলেন একটি সাধারণ \lat{mat} (চাটাই)-এ, সামনে সামান্য যব।

উমার (রাঃ) এ দৃশ্য দেখে কেঁদে ফেললেন — তাঁর \lat{tears} (কান্না) দুনিয়ার অভাবের জন্য নয়, নবীজির মহত্ত্ব ও ত্যাগের গভীর বোধ থেকে।

\medskip\hrule\medskip

\subsection*{৩. পূর্ণ ঘটনার বর্ণনা}

উমার ইবনুল খাত্তাব (রাঃ) বলেন:

\begin{quote}
\textbf{\lat{“}আমি নবীজির (সাল্লাল্লাহু আলাইহি ওয়া সাল্লাম) কাছে প্রবেশ করলাম — তিনি তখন একটি চাটাইয়ে বসা/শুয়া ছিলেন। আমি বসলাম। তখন তাঁর পরনে ছিল একটি ইযার (তহবন্দ), আর চাটাই ও তাঁর শরীরের মাঝে আর কোনো বাধা ছিল না — চাটাইয়ের দাগ তাঁর পাশে (শরীরে) বসে গেছে।} \\
\textbf{ঘরের এক কোণে দেখলাম — প্রায় এক সা' পরিমাণ এক মুঠো যব এবং কিছু গাছের পাতা। আর ঝুলানো ছিল একটি চামড়ার মশক।} \\
\textbf{এ দেখে আমার চোখ দিয়ে পানি গড়িয়ে পড়ল। নবীজি বললেন: "হে খাত্তাবের ছেলে! তোমার কান্নার কারণ কী?"} \\
\textbf{আমি বললাম: "হে আল্লাহর নবী! আমার কান্না করব না কেন? এই চাটাই আপনার পাশে দাগ বসিয়ে দিয়েছে, আর এটি আপনার সমস্ত সঞ্চয় — আমি এর ভেতরে যা দেখলাম তা ছাড়া আর কিছুই দেখছি না। অথচ কিসরা ও কায়সার ফল-ফুল ও নদীর (বিলাসে) আছে! আপনি তো আল্লাহর নবী ও তাঁর মনোনীত ব্যক্তি — আর এটাই আপনার সঞ্চয়!"} \\
\textbf{নবীজি (সা.) বললেন: "হে খাত্তাবের ছেলে! তুমি কি এতে খুশি না যে, আখেরাত আমাদের জন্য, আর দুনিয়া তাদের জন্য?"} \\
\textbf{আমি বললাম: "অবশ্যই খুশি।"\lat{”}}
\end{quote}

\medskip\hrule\medskip

\subsection*{৪. ধাপে ধাপে বিশ্লেষণ}

\textbf{ধাপ ১ — প্রবেশ:} উমার নবীজির কামরায় প্রবেশ করেন — \lat{situation} (পরিস্থিতি) জানতে।

\textbf{ধাপ ২ — দৃশ্য অবলোকন:} চাটাই, চাটাইয়ের দাগ, সামান্য \lat{barley} (যব), গাছের পাতা, একটি মশক — এটাই দুনিয়ার শ্রেষ্ঠ মানুষটির সম্পদ।

\textbf{ধাপ ৩ — উমারের কান্না:} উমার (রাঃ) সৃষ্টির সেরা মানুষটির এই \lat{simplicity} (সরলতা) দেখে কাঁদেন — অথচ দুনিয়ার দুই সম্রাট \lat{luxury}-তে মগ্ন। উমারের \lat{tears} ছিল নবীজির ত্যাগ ও আখেরাতপ্রীতি উপলব্ধির ফল।

\textbf{ধাপ ৪ — নবীজির প্রশ্ন ও জবাব:} নবীজি উমারকে \lat{comfort} (সান্ত্বনা) দিয়ে বলেন — আখেরাত আমাদের, দুনিয়া তাদের। অর্থাৎ দুনিয়ার \lat{luxury} সাময়িক, আখেরাতের সফলতা স্থায়ী। উমার এই \lat{reply} (জবাব)-এ সম্মতি জানান।

\textbf{ধাপ ৫ — চিরন্তন শিক্ষা:} এই দৃশ্য উমারের জীবনে এত গভীর \lat{impact} (প্রভাব) ফেলেছিল যে, পরে তিনি যখন খলিফা হন, তখনও সাধারণ জীবনযাপন করতেন — প্রমাণ সেটা নবীজির শিক্ষার ফল।

\medskip\hrule\medskip

\subsection*{৫. শব্দের ব্যাখ্যা}

\begin{table}[h]\centering\small
\begin{tabular}{ll}
শব্দ & অর্থ \\
\hline
\textbf{হাসীর (\arab{حصير})} & চাটাই — খেজুর পাতার শক্ত বুনন \\
\textbf{ইযার (\arab{إزار})} & তহবন্দ/লুঙ্গি — কোমরের নিচে পরার কাপড় \\
\textbf{কবযাহ (\arab{قبضة})} & মুঠো/এক মুঠো পরিমাণ \\
\textbf{সা' (\arab{صاع})} & পরিমাপের একক — প্রায় ২.৫-৩ কেজি \\
\textbf{কারায (\arab{قرظ})} & বাবলা গাছের পাতা — ট্যানিং/ঔষধে ব্যবহৃত \\
\textbf{ইহাব (\arab{إهاب})} & চামড়ার মশক — পানি রাখার পাত্র \\
\textbf{কিসরা ও কায়সার} & পারস্য ও রোমের সম্রাট — সে সময়ের দুই পরাশক্তি \\
\hline
\end{tabular}
\end{table}

\medskip\hrule\medskip

\subsection*{৬. সংশ্লিষ্ট হাদিস ও আয়াত}

\textbf{হাদিস (বুখারি ২৪৬৩):} নবীজি বলেছেন: \textbf{\lat{“}হে খাত্তাবের ছেলে! তোমার কি এতে খুশি না যে, দুনিয়া তাদের জন্য আর আখেরাত আমাদের জন্য?\lat{”}} — একই ঘটনার আরেক বর্ণনা।

\textbf{হাদিস (বুখারি ৪৯১৮):} উমার নবীজির কামরায় ঢুকে স্ত্রীদের থেকে আলাদা থাকার ঘটনা বর্ণনা করেছেন — নবীজির বিছানায় ছিল চাটাই, বালিশে খেজুরের আঁশ।

\textbf{আয়াত (সূরা আয-যুখরুফ ৪৩:৩২):}
\begin{quote}
\textbf{\lat{“}আর আপনার রবের রহমতই উত্তম — যা তারা জমা করে তার চেয়ে।\lat{”}}
\end{quote}

\textbf{আয়াত (সূরা আত-তাওবা ৯:৩৮):}
\begin{quote}
\textbf{\lat{“}তোমরা কি দুনিয়ার জীবনে খুশি হয়ে আখেরাত ছেড়ে দিয়েছ? দুনিয়ার জীবনের সামগ্রী তো আখেরাতের তুলনায় নগণ্য।\lat{”}}
\end{quote}

\textbf{আয়াত (সূরা আল-আ'লা ৮৭:১৭):}
\begin{quote}
\textbf{\lat{“}অথচ আখেরাত উত্তম ও চিরস্থায়ী।\lat{”}}
\end{quote}

\medskip\hrule\medskip

\subsection*{৭. গভীর শিক্ষা}

\textbf{১. দুনিয়ার বিলাস সাময়িক — আখেরাত স্থায়ী।}
নবীজির জীবনের \lat{simplicity} শিক্ষা দেয় — দুনিয়ার আরামের চেয়ে আখেরাতের সফলতা গুরুত্বপূর্ণ। যে এই সত্যটি অন্তরে ধারণ করে, সে দুনিয়ার \lat{arrogance} (অহংকার) ও উজবে পড়ে না।

\textbf{২. বড়ত্ব আল্লাহর নিকট — দুনিয়ার ধনে নয়।}
নবীজি দুনিয়ার সবচেয়ে বড় \lat{ruler} (শাসক) ছিলেন না — কিন্তু আল্লাহর নিকট সবচেয়ে মর্যাদাবান। দুনিয়ার বড়ত্বের \lat{arrogance} ভ্রান্ত; প্রকৃত বড়ত্ব তাকওয়ায়।

\textbf{৩. সাহাবাদের অন্তরের গভীর বোধ।}
উমার (রাঃ) নবীজির ত্যাগ দেখে কেঁদেছেন — তাদের অন্তর \lat{afterlife} (আখেরাত)-এর চিন্তায় এমনভাবে আবদ্ধ ছিল যে দুনিয়ার অভাব তাদের কাঁদায়নি, বরং নবীজির মহত্ব ও আখেরাতপ্রীতি কাঁদিয়েছিল।

\textbf{৪. ত্যাগের পুরস্কার।}
নবীজি দুনিয়া ত্যাগ করেছেন, \lat{afterlife}-এ পাবেন সর্বোচ্চ মর্যাদা। এই ত্যাগই তাঁর \lat{humility} (বিনয়)-এর প্রমাণ — আর \lat{humility}-ই \lat{arrogance} ধ্বংসের পথ।

\textbf{৫. নেতা-শাসকের জন্য সরলতা।}
উমার পরে খলিফা হয়ে নবীজির এই জীবন থেকেই \lat{lesson} (শিক্ষা) নিয়ে সাধারণ জীবনযাপন করেছেন। যারা \lat{power} (ক্ষমতা)-এ আছে, তাদের জন্য এই ঘটনা বিশেষ \lat{lesson} — \lat{power}-এর \lat{arrogance} নয়, \lat{service} (সেবা)-এর \lat{humility}।

\medskip\hrule\medskip

\subsection*{৮. জীবনে প্রয়োগের পদ্ধতি}

\begin{enumerate}
\item \textbf{নিজের সম্পদের হিসাব করুন:} আমি যা পেয়েছি তা কি আখেরাতের পথে খরচ হচ্ছে, নাকি দুনিয়ার \lat{luxury}-তে? নবীজির সামান্য সঞ্চয় আমাদের জন্য \lat{accountability} (দায়বদ্ধতা)-র দাবি রাখে।
\item \textbf{আরাম-বিলাসে অহংকার নয়:} নতুন গাড়ি, বাড়ি বা পোশাক পেলে আল্লাহর শুকরিয়া করুন, প্রদর্শন বা \lat{pride} (গর্ব) নয়।
\item \textbf{আখেরাতের চিন্তা বাড়ান:} কুরআন তিলাওয়াত, জিকির ও মৃত্যুর \lat{remembrance} (স্মরণ) — এই চিন্তা দুনিয়ামোহ দূর করে।
\item \textbf{সাহাবাদের সরলতা অনুশীলন:} নবীজি ও উমারের সরল জীবন থেকে — পরিমিত খরচ, প্রদর্শনবর্জিত জীবন — নিজে \lat{practice} (অনুশীলন) করুন।
\item \textbf{দুনিয়া ও আখেরাতের তুলনা:} দুনিয়ার কোনো বিলাসের সামনে নিজেকে প্রশ্ন করুন — এটি কি আখেরাতের চেয়ে \lat{valuable} (মূল্যবান)?
\end{enumerate}
\medskip\hrule\medskip

\subsection*{৯. রেফারেন্স}

\begin{itemize}
\item সুনানে ইবনে মাজাহ, হাদিস ৪১৫৩ (কিতাবুজ-যুহদ)
\item সহীহ বুখারি, হাদিস ৪৯১৮ ও ২৪৬৩ (দীর্ঘ বর্ণনা)
\item সহীহ মুসলিম, হাদিস ১৪৭৯ (নবীজির স্ত্রীদের থেকে আলাদা থাকার প্রসঙ্গ)
\item সূরা আত-তাওবা, আয়াত ৩৮; সূরা আয-যুখরুফ, আয়াত ৩২; সূরা আল-আ'লা, আয়াত ১৭
\end{itemize}

\end{document}
